\subsection{THE COMPLEX LOGARITHM}

Let B be the ``strip'' formed by the points \(z = x + iy\) such that \(-\pi \leqslant y < \pi\) ; the function \(e^{z}\) takes each of its values once and only once on B; in other words, \(z \mapsto e^{z}\) is a bijective continuous map of B onto \(C^{*}\) ; the image under this map of the (half-open) segment \(x = x_{0}\) , \(-\pi \leqslant y < \pi\) is the circle \(|z| = e^{x_{0}}\) ; the image of the line \(y = y_{0}\) is the (open) half-line defined by \(\operatorname{Am}(z) = y_{0} \pmod{2\pi}\) . The image under \(z \mapsto e^{z}\) of the interior B of B, that is, of the set of \(z \in C\) such that \(|\mathcal{I}(z)| < \pi\) , is the complement F of the (closed) negative real half-axis in C; if one agrees to denote by \(\operatorname{Am}(z)\) the measure of the amplitude of z which belongs to \([-\pi, \pi]\) , then the set F can be defined by the relations \(-\pi < \operatorname{Am}(z) < \pi\) . Since \(z \mapsto e^{z}\) is a strict homomorphism of C onto \(C^{*}\) the image under this map of any open subset of B (so of C) is an open set in \(C^{*}\) (so in F); in other words, the restriction of \(z \mapsto e^{z}\) to B is a homeomorphism of B onto F. We denote by \(z \mapsto \log z\) the homeomorphism of F onto B which is the inverse of the latter; for a complex number \(z \in F\) , \(\log z\) is called the principal value of the logarithm of z. If \(z = x + iy\) and \(\log z = u + iv\) then \(x + iy = e^{u + iv}\) , whence \(e^{u} = |z|\) , and since \(-\pi < v < \pi\) , we have \(v = \operatorname{Am}(z)\) . Moreover, we have \(\tan(v + \pi/2) = -x/y\) if \(y \neq 0\) ; thus we can write

\[
\left\{ \begin{array}{l l} u = \log | z | = \frac {1}{2} \log (x ^ {2} + y ^ {2}) \\ v = \frac {\pi}{2} - \operatorname{Arc} \tan \frac {x}{y} & \text { if } y > 0 \\ v = 0 & \text { if } y = 0 \\ v = - \frac {\pi}{2} - \operatorname{Arc} \tan \frac {x}{y} & \text { if } y <   0. \end{array} \right.\tag{30}
\]

It is clear that \(\log z\) is the extension to F of the function \(\log x\) defined on the positive open real half-axis \(R_{+}^{*}\) . If z, \(z'\) are two points of F such that \(zz'\) is not real negative, we have \(\log(zz') = \log z + \log z' + 2\varepsilon\pi i\) , where \(\varepsilon = +1\) , -1 or 0 depending on the values of \(\mathrm{Am}(z)\) and \(\mathrm{Am}(z')\) .

We note that at the points of the negative real half-axis the function \(\log z\) has no limit; to be precise, if x tends to \(x_{0} < 0\) and if y tends to 0 remaining \textgreater{} 0 (resp. \textless{} 0), then \(\log z\) tends to \(\log |x_{0}| + \pi i\) (resp. \(\log |x_{0}| - \pi i\) ); when z tends to 0, \(|\log z|\) increases indefinitely.

We shall see later how the theory of analytic functions allows us to extend the function \(\log z\) , and to define the complex logarithm in full generality.

Since \(\log z\) is the inverse homeomorphism of \(e^{z}\) , the formula for differentiating inverse functions (I, p.~9, prop. 6) shows that \(\log z\) is differentiable at every point \(z \in F\) , and that

\[
\mathrm{D} (\log z) = \frac {1}{e ^ {\log z}} = \frac {1}{z}\tag{31}
\]

a formula which generalizes formula (10) of III, p.~93.

